\chapter{Resumé}

Táto diplomová práca sa zaoberá modelovaním a prediktívnym riadením doskového tepelného výmenníka s využitím Koopmanovho modelu. Hlavným cieľom práce bolo identifikovať dynamiku tepelného výmenníka z experimentálnych dát, navrhnúť regulátory typu Model Predictive Control (MPC) pre dva rôzne predikčné modely a následne porovnať ich správanie v simuláciách aj pri experimentoch na reálnom zariadení. Porovnávané boli dva modelovacie prístupy: jednoduchý lineárny model prvého rádu získaný zo skokových odoziev a dátovo riadený Koopmanov model, ktorý využíva neurónovú sieť na transformáciu meraných veličín do viac rozmerného priestoru.

Výmenníky tepla patria medzi často používané zariadenia v chemickom, potravinárskom a energetickom priemysle. Ich riadenie je dôležité najmä v prípadoch, kde je potrebné presne udržiavať požadovanú teplotu výstupného média. Tepelné procesy však zvyčajne vykazujú pomalú dynamiku, obmedzenia akčných členov a nelineárne správanie. Klasické lineárne modely sú jednoduché, dobre interpretovateľné a často postačujúce v okolí pracovného bodu. Ich presnosť sa však môže zhoršiť pri väčších odchýlkach od pracovného bodu alebo pri výraznejších nelineárnych javoch. Z tohto dôvodu sa v práci skúma, či môže Koopmanov model poskytnúť lepší model a následne zlepšiť kvalitu MPC riadenia.

V teoretickej časti práce je najskôr opísaný Koopmanov operátor. Tento prístup umožňuje reprezentovať nelineárny dynamický systém pomocou lineárnej dynamiky v tzv. zdvihnutom priestore. Pôvodné stavové veličiny sú transformované pomocou zdvíhacej funkcie (lifting function) do viac rozmerného priestoru, v ktorom možno dynamiku aproximovať lineárnym modelom. Keďže presný Koopmanov operátor je nekonečno-rozmerný, v praktických aplikáciách sa používa jeho konečno-rozmerná aproximácia. V tejto práci má výsledný Koopmanov model tvar lineárneho stavového modelu v zdvihnutých súradniciach, pričom výstup systému sa získa spätným zobrazením zo zdvihnutého priestoru.

Následne je v práci predstavený prístup Deep Koopman. V tomto prípade nie je zdvíhacia funkcia zvolená ručne, ale je naučená priamo z dát pomocou neurónovej siete. Použitý model má štruktúru encoder--lineárna dynamika--decoder. Encoder mapuje meranú výstupnú teplotu do latentného stavu, lineárna časť opisuje dynamiku v latentnom priestore a decoder mapuje latentný stav späť na predikovaný výstup. V implementovanom modeli bol použitý latentný stav s dimenziou $n_z = 10$. Encoder bol tvorený plne prepojenou neurónovou sieťou s tromi skrytými vrstvami so 60 neurónmi v každej vrstve a s aktivačnými funkciami ReLU. Decoder bol lineárny, takže zodpovedal výstupnej matici Koopmanovho modelu.

Súčasťou teoretickej časti je aj opis Kalmanovho filtra, ktorý sa v práci používa na odhad stavov potrebných pre MPC. Pri lineárnom modeli slúži Kalmanov filter na odhad vnútorného stavu modelu prvého rádu. Pri Koopmanovom MPC pracuje Kalmanov filter v zdvihnutom priestore a odhaduje latentný stav, ktorý nie je priamo merateľný. Tento odhad je následne použitý ako počiatočný stav pre predikciu v MPC optimalizačnej úlohe. Použitím rovnakej štruktúry odhadu stavov pre oba regulátory sa zabezpečilo spravodlivé porovnanie jednotlivých modelov.

Ďalšou dôležitou časťou teórie je Model Predictive Control. MPC je riadiaca stratégia založená na opakovanom riešení optimalizačnej úlohy na konečnom horizonte. V každej perode vzorkovania sa pomocou predikčného modelu vypočíta optimálna postupnosť budúcich vstupov, pričom do systému sa aplikuje iba prvý vstup. Následne sa meranie aktualizuje a optimalizačná úloha sa rieši znova. Výhodou MPC je možnosť priamo zahrnúť obmedzenia na vstup aj výstup systému. V tejto práci boli pre oba regulátory použité rovnaké obmedzenia, rovnaký predikčný horizont a rovnaké váhy v účelovej funkcii, aby sa rozdiely vo výsledkoch dali pripísať najmä rozdielom modelov.

Experimentálna časť práce bola realizovaná na zariadení Armfield PCT23, ktoré predstavuje laboratórny doskový tepelný výmenník s oddeleným horúcim a studeným okruhom. Riadenou veličinou bola výstupná teplota horúceho média z výmenníka označená ako $T_4$. Manipulovanou veličinou bol riadiaci signál čerpadla horúceho okruhu Pump~F. Tento signál bol zadávaný v percentách a bol obmedzený na rozsah $20$--$100~\%$. Čerpadlo studeného okruhu Pump~W bolo počas experimentov udržiavané na konštantnej hodnote $50~\%$. Teplota horúceho média bola udržiavaná pomocným regulačným obvodom, ktorý nie je súčasťou hlavného návrhu MPC.

Identifikácia systému bola vykonaná pomocou skokových zmien vstupu. Rýchlosť čerpadla Pump~F bola menená v skokoch každých $250~\mathrm{s}$ a merala sa príslušná odozva výstupnej teploty $T_4$. Pred identifikáciou boli vstupné a výstupné dáta normalizované, aby sa pracovalo s bezrozmernými veličinami a aby sa zlepšila numerická podmienenosť výpočtov. Normalizácia bola vykonaná odčítaním empirickej strednej hodnoty a delením smerodajnou odchýlkou. Rovnaké škálovanie bolo použité aj pri Koopmanovej identifikácii a pri následnom návrhu MPC, aby boli účelová funkcia a obmedzenia vyjadrené konzistentne v normalizovaných súradniciach.

Na základe skokových odoziev bol najskôr identifikovaný lineárny model prvého rádu. Z normalizovaných skokových odoziev bola vytvorená priemerná odozva, z ktorej sa určilo zosilnenie a časová konštanta systému. Identifikované parametre boli približne $K \approx 1.17$ a $\tau \approx 68$ vzoriek. Pri vzorkovacej perióde $T_{\mathrm{s}} = 1~\mathrm{s}$ bol model prevedený do diskrétneho stavového tvaru s maticami $A = 0.985$, $B = 0.016$, $C = 1$ a $D = 0$. Tento model bol následne použitý pre lineárne MPC.

Predikčná schopnosť lineárneho modelu prvého rádu bola overená v otvorenej slučke na testovacej časti dát. Pri tomto overení bol do modelu privádzaný nameraný vstupný signál a výstup modelu nebol korigovaný meraním. Výsledkom bola čistá predikcia modelu založená len na vstupnom signáli a počiatočnom stave. Chyba bola vyhodnotená pomocou RMSE v fyzikálnych jednotkách, teda v stupňoch Celzia. Pre lineárny model bola dosiahnutá hodnota $\mathrm{RMSE}_{\mathrm{OL}} = 1.6962\,^\circ\mathrm{C}$.

Druhý predikčný model bol vytvorený pomocou Koopmanovho prístupu s použitím knižnice Neuromancer. Model bol trénovaný na normalizovaných vstupno-výstupných dátach. Použitá Deep Koopman štruktúra obsahovala encoder, lineárnu dynamiku v latentnom priestore a lineárny decoder. Tréning prebiehal na sekvenciách dĺžky $n_{\mathrm{steps}} = 50$ a s veľkosťou mini-batchov $b_s = 100$. Použitý bol optimalizátor Adam s rýchlosťou učenia $10^{-3}$, maximálny počet epoch bol 900 a tréning obsahoval early stopping na základe validačnej chyby. Optimalizačná úloha pri tréningu zahŕňala viacero členov straty: viac-krokovú výstupnú predikciu, jednokrokovú predikciu, rekonštrukčnú chybu a chybu konzistencie latentného priestoru. Výsledné chyby boli približne $\mathrm{MSE}_{\mathrm{train}} \approx 0.068$ a $\mathrm{MSE}_{\mathrm{val}} \approx 0.060$.

Koopmanov model bol tiež overený v otvorenej slučke na testovacích dátach. Dosiahnutá hodnota bola $\mathrm{RMSE}_{\mathrm{OL}} = 1.6681\,^\circ\mathrm{C}$. V porovnaní s lineárnym modelom teda došlo len k malému zlepšeniu predikcie, približne o $1.66~\%$. Výhoda Koopmanovho modelu sa teda nemusí prejaviť pri samotnej predikcii.

Pre oba modely boli následne navrhnuté MPC regulátory. Predikčný horizont bol zvolený ako $N = 20$ vzoriek. Vo úelovej funkcii boli použité váhy $Q = 10$ pre regulačnú odchýlku a $R = 1$ pre regulačný zásah. Vstup bol obmedzený na rozsah $20$--$100~\%$ a výstupná teplota na rozsah $0$--$70\,^\circ\mathrm{C}$. Optimalizačná úloha bola riešená pomocou nástroja YALMIP a riešiča \texttt{quadprog}. V oboch prípadoch bol použitý Kalmanov filter so zvolenými kovariančnými maticami tak, aby poskytoval uspokojivý odhad stavu.

Najskôr boli vykonané simulácie riadnia. Ako simulačný systém bol použitý nelineárny dátovo riadený baseline model vytvorený v prostredí Neuromancer. Simulácie boli vykonané pre viacero počiatočných teplôt $T_0 \in \{45,50,55,58,60,62,66,68\}\,^\circ\mathrm{C}$ a ich cieľom bolo overiť, či oba regulátory dokážu stabilne priviesť systém do požadovanej ustálenej oblasti. Výsledky ukázali, že oba regulátory dosahujú veľmi podobné správanie. Priemer RMSE bol $1.59\,^\circ\mathrm{C}$ pre Koopmanove MPC a $1.65\,^\circ\mathrm{C}$ pre lineárne MPC. Priemerná hodnota IAE bola 115.77 pre Koopmanove MPC a 130.41 pre lineárne MPC. Priemerné hodnoty uzavretoslučkovej účelovej funkcie boli takmer rovnaké, konkrétne 131.37 a 131.79. Simulácie teda ukázali miernu výhodu Koopmanovho MPC v sledovaní výstupu, ale bez výraznej dominancie v celkovej kvalite riadenia.

Po simuláciách boli regulátory overené experimentálne na reálnom tepelnom výmenníku. Experimenty boli spustené z rôznych počiatočných teplôt a systém bol riadený do ustáleného stavu. Požadovaná hodnota bola určená v normalizovaných súradniciach a zodpovedala cieľovej teplote približne $63.808\,^\circ\mathrm{C}$. Pre oba regulátory boli zachované rovnaké nastavenia MPC, rovnaké obmedzenia a rovnaké parametre Kalmanovho filtra.

Experimentálne výsledky ukázali, že oba regulátory dokážu stabilne riadiť reálny tepelný výmenník. Rozdiely medzi regulátormi boli výraznejšie ako v simuláciách, pretože reálne zariadenie obsahuje merací šum, hydraulické a tepelné poruchy, obmedzenia akčných členov a nelineárne tepelné javy, ktoré nie sú úplne zachytené ani jedným modelom. Koopmanove MPC dosiahol nižšiu priemernú hodnotu RMSE, konkrétne $1.011\,^\circ\mathrm{C}$ v porovnaní s $1.094\,^\circ\mathrm{C}$ pre lineárne MPC. Priemerná hodnota IAE bola 178.45 pre Koopmanove MPC a 226.13 pre lineárne MPC. Z hľadiska účelovej funkcie však lineárne MPC dosiahol mierne nižšiu hodnotu, keďže jeho regulačné zásahy boli hladšie a menej agresívne.

Osobitne bola v práci vyhodnotená prechodový interval regulácie. Dôvodom je, že mnohé experimentálne odozvy vstúpili do ustálenej oblasti ešte pred koncom celého experimentu. Prechodový interval bol definovaný ako interval od začiatku experimentu po prvý okamih, keď výstup vstúpil do pásma $\pm 0.4\,^\circ\mathrm{C}$ okolo cieľovej hodnoty a zostal v tomto pásme aspoň $N_{\mathrm{hold}} = 5$ po sebe idúcich vzoriek. Toto pásmo bolo zvolené vzhľadom na šum pozorovaný v ustálených častiach experimentov.

Pre vyhodnotenie prechodového intervalu boli použité metriky čas ustálenia v tomto intervale, IAE, ISE, ITAE a hodnota účelovej funkcie uzavretej slučky $J_{\mathrm{CL}}$ vypočítaná len na danom intervale. Keďže integračné metriky neboli normalizované dĺžkou intervalu, ich hodnoty závisia nielen od veľkosti regulačnej odchýlky, ale aj od toho, ako dlho prechodový interval trvá. Preto je potrebné tieto hodnoty interpretovať spolu s ich časom ustálenia v prechodovom intervale.

Výsledky ukázali, že Koopmanove MPC dosahuje v priemere kratší čas ustálenia. Priemerný čas ustálenia bol 53.67 vzoriek pre Koopmanove MPC a 98.67 vzoriek pre lineárne MPC. Koopmanove MPC dosiahlo aj nižšiu priemernú hodnotu IAE, konkrétne 94.59 v porovnaní s 127.20 pre lineárne MPC. Výrazný rozdiel bol pozorovaný najmä pre počiatočnú teplotu $T_0 = 62\,^\circ\mathrm{C}$, kde Koopmanove MPC dosiahlo pásmo ustálenia po 12 vzorkách, zatiaľ čo lineárne MPC splnilo podmienku ustálenia až na konci pozorovaného intervalu. Tento prípad možno považovať za výnimočný, pretože výrazne ovplyvňuje priemerné aj sumárne hodnoty. Z tohto dôvodu boli v práci uvedené aj mediány, ktoré lepšie vystihujú typické správanie regulátorov.

Rozdiely medzi regulátormi vyplývajú najmä z vlastností použitých modelov. Lineárny model prvého rádu dobre opisuje dominantnú pomalú tepelnú dynamiku systému. Je jednoduchý, interpretovateľný a vedie k hladším regulačným zásahom. Je vhodný najmä v okolí identifikovaného pracovného bodu. Koopmanov model je zložitejší, pretože pracuje v latentnom priestore s vyššou dimenziou, ale dokáže lepšie zachytiť nelineárne prechodové správanie systému. Preto Koopmanove MPC často reaguje rýchlejšie na odchýlku od cieľovej teploty, čo vedie aj k rýchlejšiemu zníženiu regulačnej odchýlky. Nevýhodou je, že regulátor potom častejšie mení vstup výraznejšie a v niektorých prípadoch sa objaví aj malý prekmit.

Záver práce je, že oba regulátory sú použiteľné na riadenie doskového tepelného výmenníka. MPC s modelom získaným zo skokových zmien predstavuje jednoduché riešenie pre prevádzku v okolí pracovného bodu a poskytuje hladšie regulačné zásahy. Koopmanove MPC je zložitejšie namodelovať, ale zlepšuje prechodové správanie systému a dosahuje nižšie integrálne kritéria kvality. Hlavným prínosom Koopmanovho prístupu teda nie je dramatické zlepšenie každej metriky, ale schopnosť dosiahnuť rýchlejšiu a presnejšiu prechodovú reguláciu pri zachovaní rovnakej štruktúry MPC. Práca tak ukazuje, že Koopmanov model môže byť vhodným rozšírením klasického lineárneho prediktívneho riadenia pre nelineárne tepelné systémy, najmä v prípadoch, kde je dôležitá kvalita prechodovej fázy.